\paragraph{补充排序与校准。}

\ReviewRankingZH{} AUROC@128采用预算内检查点的最大分数，区间通过1,000次提示分组配对bootstrap计算。该线性分类器家族内的比较，为不同特征构造提供了补充排序证据。

GRU与静态探针具有接近的预算内排序表现，AUROC@128差值为0.0050，95\%区间为[-0.0019, 0.0123]。

表~\ref{tab:review_sensitivity}比较1\%、5\%和10\%校准目标下的运行特征；已拟合模型固定，原5\%目标保留为主设定。\begin{table}[t]
\centering\footnotesize
\caption{事后探索性校准目标分析；数值为三个固定种子均值（百分比）。阈值由校准集确定，5\%为主设定；比较区间未作多重校正。}\label{tab:review_sensitivity}
\begin{tabular}{lrrr}\toprule
Method & $\alpha$ & R@128 & Test FPR \\
\midrule
Prefix TF--IDF LR & 1 & 12.9 & 0.3 \\
Current-state LR & 1 & 18.0 & 1.2 \\
Dynamic GRU & 1 & 18.5 & 1.0 \\
Prefix TF--IDF LR & 5 & 28.2 & 1.6 \\
Current-state LR & 5 & 43.7 & 5.6 \\
Dynamic GRU & 5 & 38.1 & 3.3 \\
Prefix TF--IDF LR & 10 & 43.7 & 5.5 \\
Current-state LR & 10 & 61.7 & 13.5 \\
Dynamic GRU & 10 & 56.2 & 9.0 \\
\bottomrule\end{tabular}
\end{table}


严格完成前检出要求告警检查点早于回答终点，且处于128-token预算内。以全部不安全回答为分母时，主测试静态探针和GRU的召回分别为42.8\%与35.1\%。这一时机视角补充了预算内召回曲线。
